% ============================================================
% GLOSSAIRE ET ACRONYMES
% ============================================================

\chapter*{Glossaire}
\addcontentsline{toc}{chapter}{Glossaire}
\markboth{Glossaire}{Glossaire}

\vspace{0.3cm}

\begin{longtable}{p{3.5cm} p{11cm}}

\textbf{API} & \textit{Application Programming Interface} --- Interface de programmation applicative permettant la communication entre systèmes logiciels. \\[6pt]

\textbf{ASGI} & \textit{Asynchronous Server Gateway Interface} --- Standard d'interface pour les serveurs web Python asynchrones. \\[6pt]

\textbf{BPE} & \textit{Byte Pair Encoding} --- Algorithme de tokenisation sous-mot utilisé par les LLM pour décomposer le texte en unités élémentaires. \\[6pt]

\textbf{CI/CD} & \textit{Continuous Integration / Continuous Deployment} --- Pratique d'intégration et de déploiement continus automatisés. \\[6pt]

\textbf{CORS} & \textit{Cross-Origin Resource Sharing} --- Mécanisme HTTP permettant à un serveur d'autoriser les requêtes provenant d'origines différentes. \\[6pt]

\textbf{CRUD} & \textit{Create, Read, Update, Delete} --- Les quatre opérations fondamentales de gestion des données persistantes. \\[6pt]

\textbf{CUSUM} & \textit{Cumulative Sum Control Chart} --- Algorithme de contrôle statistique de processus pour la détection de changements de moyenne. \\[6pt]

\textbf{DTO} & \textit{Data Transfer Object} --- Objet de transfert de données utilisé pour structurer les échanges entre les couches applicatives. \\[6pt]

\textbf{GHCR} & \textit{GitHub Container Registry} --- Registre de conteneurs Docker intégré à GitHub pour le stockage et la distribution d'images. \\[6pt]

\textbf{GMAO} & \textit{Gestion de la Maintenance Assistée par Ordinateur} --- Système logiciel de planification et de suivi de la maintenance industrielle. \\[6pt]

\textbf{HSTS} & \textit{HTTP Strict Transport Security} --- En-tête de sécurité HTTP imposant l'utilisation exclusive de HTTPS. \\[6pt]

\textbf{JWT} & \textit{JSON Web Token} --- Standard ouvert (RFC~7519) pour la transmission sécurisée d'informations entre parties sous forme de jeton signé. \\[6pt]

\textbf{KPI} & \textit{Key Performance Indicator} --- Indicateur clé de performance permettant de mesurer l'atteinte des objectifs. \\[6pt]

\textbf{LLM} & \textit{Large Language Model} --- Modèle de langage de grande taille, réseau de neurones profond pré-entraîné sur des corpus textuels massifs. \\[6pt]

\textbf{LLMOps} & \textit{Large Language Model Operations} --- Ensemble de pratiques pour l'exploitation des LLM en production (monitoring, gouvernance, sécurité). \\[6pt]

\textbf{MoSCoW} & Méthode de priorisation des besoins~: \textbf{M}ust have, \textbf{S}hould have, \textbf{C}ould have, \textbf{W}on't have. \\[6pt]

\textbf{ORM} & \textit{Object-Relational Mapping} --- Technique de correspondance entre les objets d'un langage de programmation et les tables d'une base de données relationnelle. \\[6pt]

\textbf{PBKDF2} & \textit{Password-Based Key Derivation Function 2} --- Fonction de dérivation de clé utilisée pour le hachage sécurisé des mots de passe. \\[6pt]

\textbf{RBAC} & \textit{Role-Based Access Control} --- Modèle de contrôle d'accès basé sur les rôles attribués aux utilisateurs. \\[6pt]

\textbf{REST} & \textit{Representational State Transfer} --- Style d'architecture logicielle pour la conception d'interfaces de programmation web. \\[6pt]

\textbf{RLHF} & \textit{Reinforcement Learning from Human Feedback} --- Technique d'entraînement des LLM par renforcement à partir de retours humains. \\[6pt]

\textbf{ROI} & \textit{Return on Investment} --- Retour sur investissement, indicateur de rentabilité d'une action ou d'un projet. \\[6pt]

\textbf{RPM} & \textit{Requests Per Minute} --- Nombre de requêtes par minute, métrique de fréquence d'appels API. \\[6pt]

\textbf{SaaS} & \textit{Software as a Service} --- Modèle de distribution logicielle où l'application est hébergée par le fournisseur et accessible via Internet. \\[6pt]

\textbf{SSR / CSR} & \textit{Server-Side Rendering / Client-Side Rendering} --- Stratégies de rendu des pages web, respectivement côté serveur et côté client. \\[6pt]

\textbf{STL} & \textit{Seasonal and Trend decomposition using Loess} --- Technique de décomposition des séries temporelles en composantes de tendance, de saisonnalité et de résidu. \\[6pt]

\textbf{TLS} & \textit{Transport Layer Security} --- Protocole de sécurité assurant le chiffrement des communications réseau. \\[6pt]

\textbf{TTL} & \textit{Time To Live} --- Durée de validité d'un élément dans un cache ou d'un enregistrement DNS. \\[6pt]

\textbf{UML} & \textit{Unified Modeling Language} --- Langage de modélisation graphique utilisé en ingénierie logicielle pour la conception de systèmes. \\[6pt]

\textbf{XSS} & \textit{Cross-Site Scripting} --- Type de vulnérabilité de sécurité web permettant l'injection de scripts malveillants. \\[6pt]

\end{longtable}
